\section{Background}
\label{sec:background}


\subsection{Congruence closure}

\emph{Congruence} ($\equiv$) requires that two terms that
that have the same function and pairwise-congruent children, must
be congruent, i.e. $f(s_1,
\ldots, s_k) \equiv f(t_1, \ldots, t_k)$ whenever $s_i \equiv t_i$ for every
$i$. \emph{Congruence closure} refers to the problem of building the smallest
equivalence relation that contains a given set of equalities and is closed under
congruence.

Typically, congruence closure is implemented using a
datastructure called an \emph{e-graph}~\cite{egraphs}, which maintains a
partition of terms into equivalence classes. More formally we consider:

\begin{definition}[]
    We define $T$ to be a set of terms over a list of function symbols $F$,
    where each $t\in T$ is either a leaf-term, for example $x, y, z$ or a
    compound term of the form $f(t_1, \ldots, t_k)$ where $f\in F$ $\text{arity}(f) = k$
    and $t_1, \ldots, t_k \in T$.
\end{definition}

We say that term $x$ is the \emph{child} of term $f(x)$ and conversely term
$f(x)$ is the \emph{parent} of term $x$.

\begin{definition}[Congruence Closure]
    Given $(F, T, E)$ where $F$ is a list of function symbols, $T$ is a set of
    terms over $F$, and $E = \{a_1 = b_1, \ldots, a_n = b_n\}$ is a set of
    equalities between terms in $T$ where $a_1, \ldots, a_n, b_1, \ldots, b_n
    \in T$ for every $i$, the congruence closure of $(F, T, E)$ is the smallest
    equivalence relation $\equiv$ over $T$ such that:
    \begin{enumerate}
        \item $a_i \equiv b_i$ for every $i$.
        \item $s \equiv s$ for every term $s$ (reflexivity).
        \item If $s \equiv t$ then $t \equiv s$ (symmetry).
        \item If $s \equiv t$ and $t \equiv u$ then $s \equiv u$ (transitivity).
        \item If $s_i \equiv t_i$ for every $i$ then $f(s_1, \ldots, s_k) \equiv
        f(t_1, \ldots, t_k)$ (congruence).     
    \end{enumerate}
\end{definition}

Congruence closure is useful in practice for SMT solvers, given a formula
$\phi = a_1 = b_1 \land \ldots \land a_n = b_n \land s \neq t$, an SMT solver
will run congruence closure on the equalities $E = \{a_i = b_i\}_{i=1}^n$ to
obtain the congruence closure $\equiv$ of $E$. Then, it will check if $s$ and
$t$ are in the same equivalence class under $\equiv$. If they are, then we 
know that $s \neq t$ is false and thus the formula $\phi$ is unsatisfiable.

Similarly, an equality saturation engine may ask: what programs are equivalent
to $P$ under a set of rewrite rules $E$. To answer this question, it can run
congruence closure on $E$ to obtain the congruence closure $\equiv$ of $E$.
Then, it may check which programs $Q$ are in the same equivalence class as $P$
under $\equiv$.

The classical algorithm of Nelson and Oppen~\cite{nelson-oppen-cc} maintains a
union-find over equivalence classes of subterms. A worklist holds classes that
have just been merged. Processing a class involves walking its \emph{parent
list}---all e-nodes that mention this class as a direct
child---re-canonicalizing each, and merging any pair that the hashcons
identifies as a duplicate. New merges produce new worklist entries; the loop
runs to a fixed point. An improvement by Downey et al.~\cite{downey} uses a 
hash-cons to detect new congruences in $O(1)$ time.

\subsection{Motivating Example}
\label{sec:bg:motivating}

To make the structure concrete, consider the formula which is solved by an SMT
solver:
%
\begin{equation*}\label{ex:formula}
\begin{aligned}
& \;\;\;\;\;\; a_0 = b_0 \land\; a_1 = b_1 \\
&\;\land\;g(f(a_0), f(a_1)) \neq g(f(b_0), f(b_1))
\end{aligned}
\end{equation*}
%

This formula consists of two equalities $a_0 = b_0$ and $a_1 = b_1$, and a
single disequality between two terms built from $f$ and $g$. The disequality
compares two applications of $g$, whose four arguments are the $f$-terms.

Congruence closure proceeds as follows. After applying the two asserted unions,
the union-find places $a_0 \equiv b_0$ and $a_1 \equiv b_1$. Then by congruence
we have that $f(a_0) \equiv f(b_0)$ and $f(a_1) \equiv f(b_1)$.

Once these equalities fire, the $g$-terms become congruent $g(f(a_0), f(a_1))
\equiv g(f(b_0), f(b_1))$. Thus, the SMT solver knows that these two terms must
be equal and thus the given formula is false.

Notice that this propagates in rounds: the initial equalities are added to the
union-find, then the four $f$-term equalities are added, and finally the
$g$-term equality is added.

\begin{definition}[Congruence Depth]
    We say that two terms are $k$-step congruent if they are congruent after $k$
    rounds of congruence closure. We say that the \emph{congruence depth} of a
    formula is the smallest $k$ such that every pair of congruent terms in the
    formula is $k$-step congruent.
\end{definition}

\begin{definition}[Congruence Width]
    For each round $r$ of congruence closure, we say that $\text{width}(r)$ is
    the number of new equalities that are added in round $r$. We define the
    total width of a congruence closure problem to be $\max_{r \in
    0\ldots\text{depth}} \text{width}(r)$, i.e.the maximum number of new
    equalities that can be added in a single round of congruence closure.
\end{definition}

In general, if we have a large number of terms that
become congruent in a single round, we can add all of these terms to the
union-find concurrently. This motivates our parallel algorithm for congruence
closure, which processes each round of congruence closure in parallel.


% : initially, the only equalities we know
% are $a_0 = b_0$ and $a_1 = b_1$, these can be added to the union-find
% concurrently. Once these are added, the four equalities between the $f$-terms
% are directly implied by these. Hence, they can be added concurrently in the next
% round. Finally, once these are added, the two $g$-terms are equal and can be
% added. Using our parallel algorithm (described in \autoref{sec:alg}), we can
% finish congruence closure in three rounds of size $2, 4, 1$ respectively. This
% is opposed to the sequential algorithm, which would process $7$ merges one at a
% time. We make these intuitions precise in Section~\ref{sec:alg}.


\subsection{Parallel algorithms}
